\uuid{vfrn}
\titre{Domaine de définition}
\auteur{bodin ; hanani ; mzari}
\organisation{exo7}
\licence{CC BY-NC-SA 4.0}
\datecreate{2020-01-20}
\niveau{L1}
\module{Analyse}
\chapitre{Continuité, limite et étude de fonctions réelles}
\sousChapitre{Autre}
\difficulte{1}
\theme{Notion de fonctions, domaine de définition}
\exo7qcmid{lille-1.continuite.003}

\contenu{
\begin{question}
\qclassification{123.00}
\qtags{motcle=domaine de définition}
Quelles sont les assertions vraies concernant le domaine de définition des fonctions suivantes ? (Rappel : le domaine de définition de $f$ est le plus grand ensemble $D_f \subset \Rr$ sur lequel $f$ est définie.)

\begin{answers}
    \good{Le domaine de définition de $\exp(\frac{1}{x^2+1})$ est $\Rr$.}

    \bad{Le domaine de définition de $\sqrt{x^2-1}$ est $[1,+\infty[$.}

    \bad{Le domaine de définition de $\frac{1}{\sqrt{(x-1)(x-3)}}$ est $]1,3[$.}

    \bad{Le domaine de définition de $\ln(x^3-8)$ est $[2,+\infty[$.}

\end{answers}
\begin{explanations}
$\frac 1x$ est définie pour $x\neq 0$, $\sqrt{x}$ est définie pour $x \ge 0$ ; $\exp x$ est définie sur $\Rr$ ; $\ln x$ seulement pour $x>0$.
\end{explanations}
\end{question}
}
